% Generated draft: requires mathematical and source review.
\documentclass[11pt]{article}
\usepackage[margin=25mm]{geometry}
\usepackage{fontspec}
\setmainfont{Latin Modern Roman}
\usepackage{amsmath,amssymb,mathtools,unicode-math}
\setmathfont{Latin Modern Math}
\usepackage{enumitem,needspace,xurl,hyperref}
\setlength{\parindent}{0pt}
\setlength{\parskip}{5pt}
\setlength{\emergencystretch}{3em}
\newcommand{\N}{\mathbb{N}}
\newcommand{\Z}{\mathbb{Z}}
\newcommand{\Q}{\mathbb{Q}}
\newcommand{\R}{\mathbb{R}}
\newcommand{\C}{\mathbb{C}}
\newcommand{\F}{\mathbb{F}}
\newcommand{\A}{\mathbb{A}}
\newcommand{\PP}{\mathbb{P}}
\newcommand{\E}{\mathbb{E}}
\newcommand{\eps}{\varepsilon}
\newcommand{\dd}{\,\mathrm{d}}
\begin{document}
\section*{UnsolvedMath ID: 30001353}
\textbf{Source problem code:} catalogue entry
\subsection{Definitions and mathematical statement}
Let \(g\) denote the genus parameter, and let \(\mathcal M_g\) denote the moduli space of curves intended by the source. The source uses \(\operatorname{Eff}(\mathcal M_g)\) for an effective divisor cone and \(\operatorname{Mov}(\mathcal M_g)\) for a movable divisor cone, and writes \(s(D)\) for the slope of a divisor \(D\). It does not supply a definition of this slope or specify the precise divisor-class spaces and cone conventions.

The requested problem is to estimate the two invariants
\[
 s(\mathcal M_g)
 =\inf_{D\in\operatorname{Eff}(\mathcal M_g)}s(D),
 \qquad
 s'(\mathcal M_g)
 =\inf_{D\in\operatorname{Mov}(\mathcal M_g)}s(D).
\]
These formulas retain the source's notation and quantification. They require clarification before they define precise numerical invariants: the source does not specify the range of \(g\), the base field, whether \(\mathcal M_g\) means the open moduli space or a compactification, whether a stack or coarse moduli space is intended, the divisor classes over which the infima range, or how to handle classes for which the slope is undefined. In particular, the displayed infima do not explicitly exclude the zero class. If boundary coefficients enter the intended slope, the compactification, boundary classes, and coefficient convention must also be supplied. No such choices are imposed in this draft.

The supplied background describes an asymptotic upper bound for the effective-cone slope and an underlying slope-limit question, but supplies neither the bound nor a precise limit assertion. These cannot be added to the displayed estimation problem without further source information.
\subsection{Short English statement}
Estimate, as functions of the genus, the infimum of divisor slopes over the effective cone and over the movable cone of the intended moduli space of curves. The supplied record leaves the slope definition, moduli-space convention, admissible divisors, and genus range unresolved. Its open-status label and literature assessment are supplied curation, not independently verified conclusions.
\subsection{Sources}
\begin{itemize}
\item Supplied local UnsolvedMath registry record, permanent ID \(30001353\), title \emph{Effective and Movable Slopes of Curve Moduli}, exact source problem code \texttt{catalogue entry}, in the Oberwolfach Reports Open Problems collection. No individual proposer is identified; the supplied proposed year is 2009.
\item Supplied report citation: \emph{Complex Algebraic Geometry (2010)}, Oberwolfach Reports, pp.\ 2483--2544, DOI \texttt{10.4171/owr/2009/44}. \url{https://doi.org/10.4171/owr/2009/44}. The record retains both the 2010 citation title and the 2009 DOI/year attribution; their relationship has not been independently checked.
\item Supplied literature reference, labeled ``Farkas--Morrison, slope counterexample'': \url{https://arxiv.org/abs/math/0209171}. The label and its association with the counterexample claim are supplied attribution and have not been independently verified.
\item Supplied literature reference, labeled ``Chen--Farkas--Morrison, slope-limit problem'': \url{https://www.claymath.org/library/proceedings/cmip018c.pdf}. The label and its association with the slope-limit question have not been independently verified.
\item Supplied literature reference, labeled ``original OWR report'': \url{https://doi.org/10.4171/owr/2009/44}.
\item Literature assessment dated exactly \texttt{2026-08-21}: the supplied curation labels the problem open, says that the record gives an asymptotic upper bound without resolving the underlying slope-limit question, and reports that later counterexamples invalidate the classical slope conjecture in some genera. This is supplied dated curation, not an independently verified literature assessment; the background itself requests independent verification of time-sensitive status and cited claims.
\item The registry assigns statement status \texttt{corrected\_verified} and reports comparison with the official Oberwolfach report, with notation, context, and question form normalized without changing the source-level claim. This is a supplied verification claim, not verification performed for this draft.
\end{itemize}
\end{document}
